\chapter[Debt contracts and financing]{Debt is a dated contract, not a balancing amount}
\label{sim:debt}

Debt issuance joins a cash transaction today to promises about future payments and creditor rights. A journal can record proceeds and a liability without describing those promises. A structural model can describe the promises without supplying the accounting measurement. Scenario simulation needs both, together with a financing policy and a model of whether counterparties accept the proposed issue. Confusing these layers can make a company appear liquid while its funding equation is economically inconsistent.

\section{Five amounts that should not share one symbol}
For a debt instrument distinguish face or promised payment $F$, gross proceeds $P$, net cash received $P-f$ after issuance costs $f$, accounting carrying amount $B$, and market value $V_D$. They may coincide in a simple par loan at inception. They need not coincide for discounted debt, transaction costs, accrued interest, modifications or subsequent market changes.

The gross contractual promise is relevant to redemption. Net cash is relevant to liquidity. Carrying amount is relevant to the balance sheet under the selected basis. Market value is relevant to valuation or repurchase. A debt-to-equity covenant may use yet another defined measure. Reusing one variable called debt for all of these quantities creates errors that can remain hidden in a par, no-fee example.

A debt record also needs currency, lender, borrower, issue and settlement dates, maturity, coupon or reset rule, day count, amortization, seniority, collateral, covenants, options and modification rights. Unknown terms remain unknown. A model can support a bounded profile without implementing every clause, but its unsupported terms must not silently default to benign values.

\section{The one-period pricing equation}
Consider a loan issued at date zero and settled at date one. Let $D$ be the date-one default indicator. If there is no default, the lender receives promised amount $F$. If default occurs, the lender receives actual recovery $R$, measured in date-one currency units. Let $r$ be a risk-free one-period rate. Under the illustrative assumption that discounting the expected payoff at $r$ is appropriate, lender break-even is
\begin{equation}
 P(1+r)=\mathbb E[(1-D)F+DR].
 \label{sim:debt-pricing}
\end{equation}
This equation is a pricing assumption for the example. Under physical probabilities, discounting at a risk-free rate requires an appropriate risk-neutrality or pricing justification. More generally, with stochastic discount factor $M$,
\begin{equation}
 P=\mathbb E[M((1-D)F+DR)].
 \label{sim:sdf-pricing}
\end{equation}
An explicitly justified pricing measure can provide an equivalent formulation. A real-world default probability is not automatically a risk-neutral probability.

Define the promised yield $y$ by $F=P(1+y)$ for positive proceeds and a single final payment. Then
\begin{equation}
 P=\frac{F}{1+y}.
 \label{sim:proceeds-face}
\end{equation}
The yield is a ratio of the promised repayment to the price, not the realized return in default. With coupons or multiple dates, a yield-to-maturity equation replaces this simple relation and may have additional conventions. The cash schedule remains the primitive description.

\section{The convention problem in the supplied note}
The user-supplied lecture note by Chak Wong adapts the risky-debt model of \citet[Section 3.6]{strebulaev2012} to residual-based deep learning \citep{user2026riskynote}. It correctly recognizes that equity value, default and debt pricing interact. Its inherited debt convention requires repair before implementation.

The retained December 2012 source's equation (3.26) raises $b'/(1+\widetilde r)$ at issuance and carries $b'$ into the next period's inherited debt. That use makes $b'$ a promised repayment. Equation (3.27), however, equates $b'(1+r)$ to an expected payoff whose nondefault leg is $b'(1+\widetilde r)$. That use makes $b'$ cash advanced. Both printed equations were inspected in the retained primary source; the mismatch is not merely a transcription error in the supplied note. The publisher's current correction status has not been cleared.

The consistent face-value convention uses $F$ for the next promised payment and $P$ for issuance cash, as in equations \eqref{sim:debt-pricing}--\eqref{sim:proceeds-face}. Substituting gives
\begin{equation}
 \frac{F(1+r)}{1+y}=\mathbb E[(1-D)F+DR].
 \label{sim:face-price}
\end{equation}
Multiplying by $1+y$ yields
\begin{equation}
 F(1+r)=(1+y)\mathbb E[(1-D)F+DR].
\end{equation}
The multiplier applies to recovery as well as to the nondefault payoff. Omitting it on recovery changes the equation when recovery is positive. Alternatively one can formulate the entire model in proceeds units, but then future liability state, repayment and recovery arguments must change consistently. Renaming one symbol in the pricing equation is insufficient.

\section{An exact positive-recovery diagnostic}
Hold face at 100, default probability at 0.1, recovery at 40 and the risk-free rate at 0.05. These are synthetic inputs chosen to expose the convention. The expected payoff is
\begin{equation}
 0.9(100)+0.1(40)=94.
\end{equation}
Thus
\begin{align}
 P&=\frac{94}{1.05}=\frac{1880}{21}\approx89.52381,\\
 y&=\frac{100}{P}-1=\frac{11}{94}\approx0.1170213.
\end{align}
The promised yield exceeds the risk-free rate because the lender does not always receive the promise. The expected payoff discounted at five per cent exactly equals the proceeds by construction.

Under the conflicting printed convention, solving $105=0.9(100)(1+y)+0.1(40)$ gives $y=11/90$. Its implied proceeds are $9000/101$. Applying the actual payoff 94 to those proceeds leaves
\begin{equation}
 94-1.05\left(\frac{9000}{101}\right)=\frac{44}{101}\approx0.43564.
\end{equation}
This is a lender break-even residual measured in date-one cash. It is not a rounding error. Zero recovery or no default can conceal the mismatch, so a no-default smoke test is inadequate. Positive recovery is the discriminating case.

\section{Recovery needs its own definition}
Recovery can mean a fraction of face, a fraction of market value before default, a share of realized collateral, or a cash distribution after costs and delays. These meanings are not interchangeable. If recovery is $R=\rho F$, the one-period expected payoff is $F[1-p+p\rho]$ when default probability is $p$. If $R$ is an absolute asset distribution, scaling face does not generally scale recovery.

For delayed recovery at date two, discounting it as a date-one payment is wrong. A stochastic discount-factor formulation must use the payment's actual date. Recovery uncertainty can be correlated with default and aggregate conditions; replacing it by an unconditional mean can misprice the claim.

Multiple creditors require an allocation rule. Suppose distributable assets after costs are 70, senior claims are 60 and junior claims are 50. Under an assumed strict priority waterfall without other rights, senior recovery is 60 and junior recovery 10. Paying both creditors 70 per cent would distribute 77 and violate the available asset constraint. The example illustrates conservation and priority, not a legal conclusion for a particular jurisdiction. Actual enforceability, security and restructuring rules need admitted authority.

\section{Default is a transition, not deletion}
An inability or decision not to pay can lead to arrears, cure periods, waiver, acceleration, restructuring, enforcement or liquidation. The contract and legal process determine which transitions are available. A simulator that sets debt and equity to zero at a default flag loses unpaid claims, assets, recovery timing and operating consequences.

A useful state includes scheduled obligation, amount paid, arrears, default status, acceleration status and recovery process. A missed payment can create an arrears balance even if operations continue. A covenant breach can restrict dividends without immediate acceleration. A waiver can restore eligibility subject to fees and revised terms.

Limited liability concerns shareholders' obligations and the value of their claim under the model. It does not require book equity to remain nonnegative. Negative book equity, negative market equity before applying limited liability, cash shortage and legal default are different concepts. Their relationships must be modeled rather than inferred from a single sign test.

\section{Refinancing creates a new problem}
A maturing debt of 100 cannot be repaid by writing a new debt balance of 100 unless a financing transaction supplies the cash in time. Suppose new face 100 sells for 90 and issuance fees are 2. Net proceeds are 88, leaving a shortfall of 12 before considering cash reserves. Treating face as cash hides that shortfall.

If the company requires net proceeds of 100 and the price is 0.9 per unit face with a fee of 2, required face is $(100+2)/0.9=113.333\ldots$. Increasing face may itself worsen price or violate capacity. The funding calculation is then coupled to the pricing and eligibility problem.

A committed facility may provide capacity but still have draw conditions, currency limits, borrowing-base constraints or expiry dates. An uncommitted facility is not guaranteed funding. A refinancing policy should distinguish request, offer, acceptance and settlement. Failure at any stage leaves a residual funding need and possible distress transitions.

\section{Cash, debt and net debt}
Some structural models use a signed debt variable, with negative debt interpreted as cash. This is economical in a stylized model, but cannot automatically represent a company with cash and debt simultaneously. Restricted cash may not repay a particular obligation. Cash in one subsidiary or currency may not be available to another. Gross debt can affect covenants and refinancing even when net debt is small.

Consider cash 100 and debt 100. Net debt is zero. If cash is unrestricted in the borrowing entity and repayment is allowed without cost, netting may be adequate for a particular valuation question. If cash is pledged to another creditor or held in a subsidiary subject to transfer restrictions, netting destroys liquidity information. The same net balance describes different payment capacity.

AccDSL should therefore preserve gross instruments and cash accounts in the general representation. A net-debt reduction is a query-specific abstraction requiring explicit conditions. It should not become the universal state merely because a numerical paper uses it successfully in a simpler economic model.

\section{Coupons, accrual and effective interest}
Contractual coupon determines a payment schedule. Accounting interest expense can differ because of discount, premium, fees or measurement rules. For a simple illustrative amortized-cost schedule with effective rate $i$, carrying amount evolves as
\begin{equation}
 B_{t+1}=B_t+iB_t-\text{cash coupon}_{t+1}-\text{principal paid}_{t+1},
\end{equation}
subject to the selected accounting treatment and any additional events. The effective-interest term is not automatically a cash payment.

Suppose a zero-coupon promise of 100 is issued for 90 for one year. With no fees, an effective rate $100/90-1$ accretes carrying amount from 90 to 100. Cash is received at issuance and paid at maturity. Charging 10 as both accretion and a coupon payment would double count the financing cost.

Conversely, a floating coupon can reset while an accounting measurement follows its own applicable rules. A simulation should generate contractual cash first from terms, then calculate accounting consequences under an admitted basis. A generic rate shock on interest expense cannot substitute for these two mechanisms in every instrument.

\section{Issuance costs and tax timing}
A fee can be paid to an intermediary at issuance, deducted from proceeds, financed separately or incurred conditionally. Its cash timing should be explicit. The accounting allocation between expense and carrying amount depends on the applicable standard and facts. A model cannot infer that allocation merely from the fact that cash was paid.

The retained structural source and supplied lecture include a simplifying treatment that credits the present value of a future interest tax benefit at issuance. That can be part of an economic valuation convention. It must not be posted as an actual contemporaneous tax receipt without a separate tax mechanism. A present value is not a cash settlement.

Tax losses, deduction limits, jurisdiction, timing and the availability of taxable income can all affect realization. A scenario that assumes every interest expense produces immediate cash tax savings may overstate liquidity in distress. The proposed design keeps tax cash payments, tax expense, deferred tax and economic tax-benefit valuation separate, then links them under explicit rules.

\section{Options and amendments}
A callable bond gives the issuer an exercise right; a puttable bond gives a holder a right. Their exercise can depend on rates, liquidity, transaction costs and information. Replacing the schedule with the contractual final maturity ignores these options. Exercising automatically whenever it appears favorable assumes a policy and ignores possible constraints.

An amendment can change coupon, maturity, collateral or covenants. It needs an effective date, consent requirements, fees and accounting treatment. A revised schedule should supersede future obligations consistently while retaining the original agreement and past payments. Applying both old and new coupons after amendment creates duplicate obligations; erasing the old contract destroys provenance.

Where a clause is unknown, the family can contain alternative exercise rights or amendment paths. Those alternatives should be activated only by their premises. A model should not give the borrower an option that belongs to the lender, nor assume the lender consents to a concession because it improves borrower liquidity.

\section{Pricing and policy interact}
A debt price can depend on future investment, payouts and refinancing because these affect repayment. The issuer's policy can depend on the price offered. Solving the two separately while holding inconsistent policies fixed can produce a price that is correct for a different company behavior.

The structural-finance survey \citet[Sections 3.4--3.6]{strebulaev2012} is useful for this interaction. The newer working paper \citet[Sections 2.1.3--2.1.5 and 3]{duarte2026} also solves bond pricing and equity value jointly under its own timing, recovery and tax assumptions. Those assumptions differ from the lecture's model; it is not a drop-in correction.

A general AccDSL financing service can accept a declared price or offer as an input. That supports conditional simulation without claiming equilibrium. An optional equilibrium service may solve pricing and policy together, with separate numerical evidence. The interface should identify which service supplied the offer and what assumptions support it.

\section{Credit supply and rationing}
A fair price does not imply an unlimited quantity is available. Lenders can face capital, liquidity, concentration or mandate constraints. A borrower may be unable to issue enough securities at any admissible terms. A model that solves only lender break-even can miss rationing.

Bank stress and balance-sheet models such as \citet{halaj2016,behn2019} provide examples where financing and adjustment depend on constraints. Their institutional settings should not be transferred mechanically to a corporate issuer. The general lesson is that a financing transaction needs a counterparty and feasible quantity, not only an interest rate.

An offer curve can therefore specify price, maximum quantity, conditions and expiry. A policy chooses among feasible offers or records no acceptable offer. If the curve is uncertain, results remain conditional on alternative credit-supply assumptions. A scenario can deliberately freeze refinancing, but that restriction should be visible as a stress premise.


\section{What dynamic financing models establish}
For longer-lived debt, lender receipts depend on future issuance, investment,
priority changes and default choices. A candidate policy $\pi$ produces a
price $q^\pi$. If management optimizes while holding an incompatible price
fixed, the resulting policy may not be the one lenders priced. A simultaneous
solution must make the chosen policy optimal under its price and make the
price value that policy's promised and default-state payments. A constrained
contract-price input is a useful alternative for mechanical stress testing,
but its interpretation is conditional on that price and availability.

Benzoni and coauthors make this distinction explicit in a dynamic model with
fixed issuance costs \cite[Section~2, Propositions~1--2, Appendix~A]{litclosebenzoni}.
Their debt quantity is face value, its price is per unit of face, and
shareholders choose default, issuance or inaction. Value matching and smooth
pasting generate a candidate barrier policy; off-policy deviations and a
supermartingale condition are required for equilibrium verification. The
retained model assumes zero default recovery and risk-neutral agents. It
therefore illustrates the policy-price consistency problem, but is not the
positive-recovery solution required by the preceding example.

The alternative managerial model of Lambrecht and Myers places investment,
borrowing, payout and rents in a dated sources-and-uses constraint and imposes
a no-Ponzi restriction \cite[Sections~2--4]{litcloselambrecht}. Its safe-debt
solutions depend strongly on managerial preferences and governance. In its
exponential-utility example, the liquidation covenant protects lenders under
the stated diffusion dynamics. A jump across the boundary, delayed liquidation
or costly asset disposal would require a different recovery analysis. A
closed-form policy derived under those restrictions is a competitor to
consider, not an identified default policy for every issuer.

Dewatripont, Legros and Matthews provide a primary contracting comparison
\cite[Sections~2--6 and Appendix~B]{litclosedewat}. Their environment has observable
but noncontractible effort, limited liability, monotone investor receipts and
efficient renegotiation. Under entrepreneur-offer renegotiation, a debt
contract can dominate a more general contract in the stated pure-strategy
class. Their appendix also supplies cases in which changed bargaining power,
third parties or mixed messages change that conclusion. Even their investor
conversion option need not be an ordinary equity security, as footnote 16
explains. Contract labels cannot replace the exact contingent payments and
exercise rights in the company model.


\section{Equal balances after unequal payment histories}
\label{sim:settlement-remedies}

A borrower can recover from a cash shortage before the reporting date. Its
closing statements then need not reveal that it failed to pay on time. To
isolate the timing, consider a synthetic company with cash 60, inventory 100,
debt carrying amount 80 and book equity 80. The contractual promise is 100,
payable on 30 June 2027. Assume a zero coupon, no fees or taxes, and one
remaining measurement interval that accretes the discount of 20 at maturity.
This deliberately simple interval makes the cash and contract comparisons
exact; a general effective-interest schedule needs its own measurement rules.

The company issues 40 new shares for one unit of cash each. In the first
history, the subscription settles on 29 June. Cash rises from 60 to 100,
the discount accretes on 30 June, and the borrower pays the complete promise
that day. In the second history, the shares settle on 2 July. The borrower
can pay only 60 at maturity. On 1 July an explicit observation records the
unpaid 40; the new cash arrives on 2 July and pays the remaining claim on
3 July. No overdraft, emergency loan or retroactive receipt funds the earlier
payment.

The following comparison follows the money and claim in the second history.
The liability column is positive for presentation, although journal credits
use the implementation's signed convention.
\begin{center}
\begin{tabular}{lrrr}
\toprule
After the dated event & Cash & Claim unpaid & Historical miss\\
\midrule
30 June: accretion &60&100&0\\
30 June: partial payment &0&40&0\\
1 July: missed-payment observation &0&40&40\\
2 July: share proceeds &40&40&40\\
3 July: final payment &0&0&40\\
\bottomrule
\end{tabular}
\end{center}
The zero in the first two historical-miss cells records that the observation
has not yet been made. It is not a finding that a payment was timely. The
observation changes contract memory without changing cash or posting another
expense. The final payment clears the current claim; it leaves the history
of the earlier shortfall intact.

Both histories end with the same financial statements under the supplied
treatment. Inventory of 100 is the only remaining asset, the debt is zero,
and equity is opening equity plus the subscription less the discount expense:
\begin{equation}
 A=100,\qquad L=0,\qquad E=80+40-20=100.
\end{equation}
Profit for the illustrated movements is $-20$. Net financing cash flow is
$40-100=-60$ in either history. The late path divides the repayment into
$60+40$, which changes its timing but not its total. Even agreement across
the closing balance sheet, income and cash-flow statements does not answer
whether the maturity was paid on time. A simulation query about missed
payments must retain the distinguishing history.

The difference also explains why balancing the journal is insufficient. Share
proceeds can debit cash and credit capital correctly while a faulty method
fails to increase the recorded shares. A missed-payment observation can omit
its history update without unbalancing any account. In the checked reference,
both omissions pass generic accounting execution and fail the independent
complete-state comparison. Explicit financing conditions subsequently reject
these proposed transitions before accepting their successors.

Those conditions have two roles. Snapshot conditions constrain current
quantities: claims and arrears are nonnegative, arrears cannot exceed the
remaining claim, and under this matured zero-coupon convention the carrying
liability equals that claim. Transition conditions constrain how an event
changes them. Payment of $p$ gives remaining claim $q'=q-p$, reduces cash by
$p$, and leaves the historical missed amount unchanged. A share subscription
must increase both the monetary capital account and the separate share
quantity under the supplied issue terms. An internally plausible rewrite of
historical miss from 40 to 50 can satisfy snapshot inequalities but violates
the transition condition. This is why both kinds of condition are needed.

Opening history creates another distinction. A supplied fiscal opening must
satisfy the declared period-start, temporary-balance and flow restrictions;
an intermediate state containing this year's payments is not such an opening.
Importing a consistent snapshot does not establish that it includes every
earlier missed payment. The selected opening method therefore compares it
with a separately supplied history premise and preserves unknown completeness.
Actual source evidence is still needed to establish that premise.

\section{A missed payment and the available remedies}
The amount unpaid does not by itself identify the creditor's remedy. Apple's
2024 annual report provides a concrete historical description of why the
distinction matters. Exhibit 4.1 expressly calls its description incomplete
and qualifies it by the 2013 and 2018 indentures; its issuer references concern
Apple Inc., rather than automatically every subsidiary
\citep[Exhibit 4.1, printed p.~3]{apple2024}. We use the summary to identify
conditions a simulation must represent. It does not establish an actual
issuer default, complete operative terms, or the applicability of those terms
to the synthetic company.

The described interest-payment event requires default continuing for 30 days
after the installment becomes due. Principal and premium use a different
condition: nonpayment when due and payable, including maturity, redemption
or declaration. Importing the interest grace period into principal repayment
would therefore apply the wrong described condition. The other-covenant route
has its own exclusions, a continuing breach and 90 days after qualifying
written notice. Holder notice uses at least 25\% of outstanding principal
under the 2013 indenture and at least 33\% under the 2018 indenture; the
trustee has a separately described notice route
\citep[Exhibit 4.1, printed p.~9]{apple2024}.

The thresholds concern the relevant outstanding series, not a fraction of
consolidated debt or a count of creditors. Series mapping, denominator,
delivery to the required recipients and validity of notice are substantive
inputs. No notice and no evidence about notice lead to different answers.
Likewise, day-level data at exactly the 30- or 90-day boundary may leave
event order unresolved. The conditional reference retains that uncertainty;
it does not assert that the complete contract has no answer at the boundary.

The summary next distinguishes declaration from automatic acceleration.
Certain continuing nonbankruptcy events permit a declaration by the specified
trustee or holder route. Certain bankruptcy events make the obligations due
without that declaration. A negative equity balance establishes neither the
specified bankruptcy event nor the required category. Moreover, an event for
one series does not necessarily constitute an event for another. Neither
automatic contagion nor universal independence follows
\citep[Exhibit 4.1, printed pp.~9--10]{apple2024}.

Cure and rescission must also remain separate. The described rescission of a
declaration requires holder action and specified deposits and cure or waiver
conditions, with an exception concerning amounts due solely because of
acceleration. Payment can remove current arrears without by itself proving
that an earlier acceleration has been rescinded. In the selected synthetic
continuation, payment changes the economic claim, reassessment changes the
current clause finding, and a separate supported rescission assessment changes
the recorded acceleration state. The historical miss remains 40 throughout.

Consent cannot be supplied by a numerical convenience either. The summary
requires each affected holder's consent for specified protected changes,
including changes to stated maturity and reductions of principal or interest.
Its general majority-waiver provision excludes payment defaults and certain
individually protected provisions. A waiver concerns the specified past
default; it does not license a future one
\citep[Exhibit 4.1, printed pp.~11--12]{apple2024}. A current contractual default
may cease under a valid waiver while the economic record still shows that
payment arrived late. Keeping those two statements is consistent.

The surrounding qualifications matter economically. Individual payment
enforcement differs from collective proceedings; legal discharge differs
from covenant defeasance. The description warns that assets sufficient to
fund scheduled maturity may be insufficient after acceleration
\citep[Exhibit 4.1, printed pp.~10--13]{apple2024}. A simulator must therefore
retain the relevant rights, legal owner, timing and usable funding, even
when a consolidated closing statement appears fully funded.

The typed conditional implementation translates the retained examples and
records selected findings through the common monetary and memory engine.
An unresolved finding produces no successor state. Recording a known finding
changes no cash, carrying amount, shares or historical missed amount. This
establishes a limited representation and execution result. Multiple concurrent
clauses, the full instruments, actual notices and amendments, and the economic
consequences of acceleration still require their own admitted methods.


\section{Implementation proposal}
The debt module should maintain immutable contract versions, a dated obligation schedule, current economic state and separate accounting measurements. Events include issue, draw, reset, accrue, coupon, amortize, redeem, amend, breach, cure, accelerate, restructure and recover. Each event has preconditions, amounts, dates and successor state.

Existing repository methods already cover bounded issue, accrual, coupon, redemption and conversion calculations, and newer declared methods preserve timing and debt memory. They are useful foundations. They do not by themselves establish a general pricing model, market capacity, legal default process or complete instrument admission. The extension should compose with their checked boundaries rather than describe those capabilities as absent.

The smallest discriminating financing fixture should include proceeds different from face, positive recovery, a fee paid at issuance, a maturity shortfall and an unavailable refinancing offer. A second fixture should add two creditors with priority and delayed recovery. A third should test unknown rate terms and a covenant boundary. Passing these checks establishes the specified mechanics; empirical pricing validity and legal applicability require separate evidence.
